\centerline{\bf \Huge По касательной}

\begin{flushright}
        \scriptsize{Вся красота мира легко может поместиться в голове одного человека.\\ Хаяо Миядзаки}
\end{flushright}

Основные свойства касательной:

\begin{enumerate}
    \item[(a)] Радиус окружности, проведенный в точку касания, перпендикулярен касательной. 
    
    \item[(b)] Отрезки касательных, проведенных к окружности из одной точки, равны.
    
    \item[(c)] \textbf{Теорема об угле между касательной и хордой.} Угол между касательной и хордой, проведёнными через общую точку на окружности, равен половине дуги, заключённой между сторонами этого угла.
\end{enumerate}


\problem{1}
Сколько различных общих касательных можно построить к двум окружностям? Рассмотрите все случаи взаимного расположения окружностей.

\problem{2}
Вневписанная окружность $\omega$ треугольника $ABC$ касается продолжений сторон $BA$ и $BC$ в точках $P$ и $Q$ докажите, что $P_{ABC} = BP + BQ$

\problem{3}
В остроугольном треугольнике ABC проведены высоты $BB_1$ и $CC_1$. Через точку $A$ проводится прямая $l$, параллельная прямой $B_1C_1$. Докажите, что эта прямая касается окружности описанной около треугольника $ABC$.

%\problem{2}В равнобедренный треугольник с основанием, равным $a$, вписана окружность и к ней проведены три касательные так, что они отсекают от данного треугольника три маленьких треугольника, сумма периметров которых равна $b$. Найдите боковую сторону данного треугольника.

\problem{4}
Касательная к описанной окружности треугольника $ABC$, проведенная через точку $A$, пересекает прямую $BC$ в точке $X$. Биссектриса угла $A$ пересекает сторону $BC$ в точке $L$. Докажите, что треугольник $ALX$ равнобедренный. Сформулируйте и докажите аналогичное утверждение для биссектрисы внешнего угла $A$.

%\problem{3}Хорда большей из двух концентрических окружностей касается меньшей. Докажите, что точка касания делит эту хорду пополам.

%\problem{4}Докажите, что если касатальная и секущая проведены к окружности из одной точки, лежащей вне окружности, то квадрат расстояния до точки касания равен произведению длины всей секущей на ее внешнюю часть.

\problem{5}
Диагонали вписанного четырёхугольника $ABCD$ пересекаются в точке $K$. Докажите, что касательная в точке $K$ к описанной окружности треугольника $ABK$ параллельна $CD$.

\problem{6}
Около треугольника $ABC$ описана окружность с центром $O$. Вторая окружность, проходящая через точки $A, B, O,$ касается прямой $AC$ в точке $A$. Докажите, что $AB = AC.$

%\problem{6}Две окружности пересекаются в точках $A$ и $B$. Через точку $K$ меньшей окружности проведены прямые $KA$ и $KB$, вторично пересекающие большую окружность в точках $P$ и $Q$ соответственно. Докажите, что хорда $PQ$ большой окружности перпендикулярна диаметру KM меньшей окружности.

%\problem{9}\textbf{Лемма Архимеда.} Окружность $\omega$ касается хорды $MN$ окружности $\Sigma$ в точке B, а самой окружности $\Sigma$ в точке $A$. Докажите, что $AB$ является биссектрисой угла $MAN.$

%\problem{10}Две окружности пересекаются в точках $A$ и $B$. К этим окружностям проведена общая касательная, которая касается окружностей в точках $C$ и $D$. Докажите, что прямая $AB$ делит отрезок $CD$ пополам.
